\chapter{Skin Cancer Screening}
\section*{What Is This Test?} Skin-cancer screening is a visual examination for suspicious lesions and is not a single blood test.\section*{Alternative Names} Skin examination; melanoma screening; skin-lesion check.\section*{Principle} Inspection and dermoscopy identify asymmetry, border, color, diameter, evolution, and other concerning features.\section*{Why This Test Is Ordered} It is used for changing lesions, personal or family history, high sun exposure, many moles, or immunosuppression.\section*{Primary vs Secondary Test} Examination is primary; biopsy provides diagnosis for a suspicious lesion.\section*{How This Test Is Performed} The clinician examines exposed and sometimes full-body skin, scalp, nails, and mucosa, often using dermoscopy.\section*{What Preparation Is Needed} Remove makeup and nail polish when possible and bring prior photographs or dermatology records.\section*{Understanding Results} Most lesions are benign, but change or concerning morphology warrants biopsy rather than reassurance from appearance alone.\section*{Follow-up Tests} Dermoscopy monitoring, biopsy, pathology, excision, or oncology assessment may follow.\section*{References} \begin{enumerate}\item MedlinePlus. Skin Cancer Screening.\item American Academy of Dermatology. Skin cancer screening resources.\end{enumerate}\clearpage\section*{Understanding Results and Follow-up}
The laboratory report should identify the specimen, method, units, reference interval or decision limit, and whether the result is positive, negative, abnormal, or inconclusive. Reference intervals vary with age, sex, pregnancy, specimen type, assay, and laboratory; a result outside a range is not automatically a diagnosis, and a result within a range does not always exclude disease. Discuss unexpected results with the ordering clinician, who can decide whether repeat or confirmatory testing, treatment, or referral is appropriate. Do not change medicines or delay urgent care based on an isolated result.
\section*{Regulatory Status and Authorization Date}This chapter describes a test category, not a specific commercial product. FDA, CDSCO, EU IVDR, and MHRA approval or registration dates therefore cannot be assigned without the assay manufacturer, product name, instrument, and intended use. For laboratory-developed tests, an individual agency approval date may not exist. See the Preface for the interpretation of regulatory status.
\clearpage
